\subsection{One-Time Unforgeability}
\label{app:signatures}

The experiment uses the signature algorithms of Section~\ref{sec:signatures}.
\begin{definition}[One-time unforgeability]
\label{def:one-time-signatures}
Give an efficient adversary these commands:
\begin{description}
  \item[$\GetVerificationKey()$.]
    Generate $(\sigvk,\sigsk)$ and return $\sigvk$.
  \item[$\SignMessage(m)$.]
    Record $m$ and return $\Sign(\sigsk,m)$.
  \item[$\Forge(m^*,\sig^*)$.]
    Terminate.  The adversary wins if
    $\Verify(\sigvk,m^*,\sig^*)=\top$ and $m^*$ was not the
    message submitted to $\SignMessage$.
\end{description}
The first two commands are accepted at most once each, in that order.
The adversary may attempt its forgery after receiving the verification
key, whether or not it requested a signature.  Malformed commands are
rejected without changing state.  The scheme is one-time unforgeable
if every efficient adversary's winning probability is negligible.
\end{definition}

In Figure~\ref{fig:signed-encrypted-inputs}, the message
$q=\InputRecord(\ctx,\xpub,c)$ contains the sampling description and
hence the provider's verification key. Applying Lamport's scheme to a
collision-resistant digest of $q$ fixes the verification-key and signature
lengths independently of the message length.
